\chapter{Metodología y planificación del trabajo}

\section{Enfoque metodológico}

Este proyecto se ha desarrollado siguiendo un modelo \textbf{incremental-iterativo} que combina las ventajas de una definición temprana de requisitos con la flexibilidad necesaria para refinar el resultado en ciclos sucesivos:

\begin{itemize}
    \item \textbf{Incrementalidad}: el software se construye por módulos funcionales independientes, cada uno de los cuales añade nuevas prestaciones al sistema global.
    
    \item \textbf{Iteración}: cada módulo se somete a múltiples rondas de desarrollo, validación y refactorización. Esto garantiza la detección rápida de errores y su corrección.
    
    \item \textbf{Adaptación evolutiva}: la metodología se adaptó según los fallos encontrados, especialmente en los múltiples intentos de integración entre módulos.
    
    \item \textbf{Validación experimental}: en todas las etapas se emplean casos de prueba reales (dispositivos sintéticos, controladores físicos, cascos EEG) para contrastar el funcionamiento.
\end{itemize}

La metodología evolucionó durante el desarrollo para responder a desafíos técnicos específicos, como la migración de Julia a Python por eficiencia de desarrollo, múltiples fallos de integración directa, y la eventual adopción de arquitecturas API como solución robusta.

\section{Fases y actividades}

\begin{table}[hp!]
\centering
\resizebox{\textwidth}{!}{
\begin{tabular}{c|l|c}
\rowcolor{udcpink!25}
\textbf{Fase} & \textbf{Actividad} & \textbf{Período} \\
\hline
0 & Investigación de tecnologías BCI y estado del arte & Marzo 2025 \\
1.1 & Desarrollo inicial script captura de inputs & 12/04-13/04 \\
1.2 & Intento de migración a Julia (fallido) & 13/04 \\
1.3 & Desarrollo UI para módulo de inputs (tkinter → PyQt5) & 14/04-15/04 \\
2.1 & Desarrollo inicial módulo procesamiento EEG & 14/04-16/04 \\
2.2 & Expansión funcionalidades EEG (brain map, ayuda, CSV) & 16/04-17/04 \\
3 & \textbf{Primer intento integración} (monolítica con LSL) & 17/04-19/04 \\
4 & Refactorización SOLID de ambos módulos & 20/04-21/04 \\
5 & \textbf{Segundo intento integración} (post-refactorización) & 22/04-30/04 \\
6 & \textit{Pausa por otros compromisos académicos} & 01/05-13/05 \\
7.1 & Desarrollo APIs de acceso a datos & 14/05-18/05 \\
7.2 & Explorador web individual por módulo & 18/05-21/05 \\
7.3 & Refactorización a explorador web unificado & 21/05-22/05 \\
8.1 & APIs de control módulo inputs & 23/05-06/06 \\
8.2 & Implementación controlador PS4 & 04/06-06/06 \\
8.3 & Tipado masivo y pre-commit & 06/06-08/06 \\
9.1 & CLI inputs (con problemas headless identificados) & 09/06-12/06 \\
9.2 & Optimizaciones PS4 y arquitectura por eventos & 10/06-12/06 \\
10.1 & API control módulo EEG (básico) & 13/06-14/06 \\
10.2 & Integración MNE y centralización DataAccessService & 15/06-17/06 \\
10.3 & Funcionalidades avanzadas (ICA, wavelets, filtros) & 17/06-18/06 \\
10.4 & Detección automática canales y streamers & 19/06-20/06 \\
10.5 & Migración lógica filtros y corrección grabación & 20/06-21/06 \\
10.6 & Finalización API control EEG y CLI & 21/06-22/06 \\
11 & \textbf{Módulo integración final} con LSL (pausado) & 22/06-23/06 \\
12 & Documentación técnica y memoria & 22/06-23/06 \\
\end{tabular}
}
\caption{Cronograma detallado del proyecto (marzo-junio 2025)}
\label{tab:cronograma}
\end{table}

\subsection{Detalles por fase}

\textbf{Fase 0: Investigación previa (marzo 2025)}
\begin{itemize}
    \item Análisis de librerías BCI y lenguajes soportados
    \item Estudio de cascos EEG comerciales y características
    \item Investigación sistema 10/20 y tipos de artefactos
    \item Evaluación inicial Julia vs Python
\end{itemize}

\textbf{Fase 1: Desarrollo módulo inputs (12-15 abril)}
\begin{itemize}
    \item Script básico con pynput: teclado → filtros → ratón complejo → clicks
    \item Intento fallido de migración a Julia por debugging complejo
    \item UI inicial en tkinter (fallida por rendimiento)
    \item Experimentación con PyQt5 via módulo EEG
    \item Migración definitiva a PyQt5 y refinamiento visual
\end{itemize}

\textbf{Fase 2: Desarrollo módulo EEG (14-17 abril)}
\begin{itemize}
    \item Experimento inicial con BrainFlow y placa sintética
    \item Soporte todos los cascos BrainFlow y calidad de señal
    \item Brain map (distribución electrodos sistema 10/20)
    \item Sistema de ayuda contextual por falta conocimiento EEG
    \item Logger CSV con métricas mentales y marcadores
\end{itemize}

\textbf{Fase 3: Primer intento integración (17-19 abril)}
\begin{itemize}
    \item Integración visual exitosa y coordinación grabadores
    \item Soporte CSV en módulo inputs para unificar logs
    \item \textbf{Fallo}: LSL con arquitectura monolítica por tamaño archivos
\end{itemize}

\textbf{Fase 4: Refactorización SOLID (20-21 abril)}
\begin{itemize}
    \item Reestructuración completa módulo inputs (día 20)
    \item Reestructuración completa módulo EEG (día 21)
    \item Archivos granulares, servicios independientes, patrones diseño
\end{itemize}

\textbf{Fase 5: Segundo intento integración (22-30 abril)}
\begin{itemize}
    \item Nuevo módulo integrador con arquitectura refactorizada
    \item \textbf{Segundo fallo}: persistencia problemas integración
    \item Mejoras interfaz y widgets ayuda formatos salida
\end{itemize}

\textbf{Fase 7: Arquitectura API (14-22 mayo)}
\begin{itemize}
    \item APIs acceso datos (módulo inputs primero, modulo EEG después)
    \item Explorador web FastAPI para pruebas API
    \item Visualizadores tiempo real WebSocket independientes
    \item Refactorización a explorador web genérico unificado
    \item Corrección problema polling concurrente
\end{itemize}

\textbf{Fase 8: APIs control (23 mayo - 8 junio)}
\begin{itemize}
    \item API control módulo inputs con sincronización GUI
    \item Plantilla inicial API control EEG
    \item Implementación controlador PS4 con SDL
    \item Unificación servidores web (datos + control)
    \item Tipado masivo y pre-commit
\end{itemize}

\textbf{Fase 9: CLIs y optimizaciones (9-12 junio)}
\begin{itemize}
    \item CLI módulo inputs basado en APIs
    \item \textbf{Identificación problema}: modo headless no funcional (dependencias Qt)
    \item PS4 no funcional en CLI (timers Qt)
    \item Migración a arquitectura dirigida por eventos
    \item Optimizaciones masivas: batching, workers, buffer circular, orjson
\end{itemize}

\textbf{Fase 10: Funcionalidades avanzadas EEG (13-22 junio)}
\begin{itemize}
    \item API control EEG: conexión placa, streams, parámetros
    \item Integración MNE-Python, centralización DataAccessService
    \item ICA, transformadas wavelet, filtros avanzados
    \item Detección automática canales, presets dinámicos
    \item Streamers archivo/TCP/UDP con restricciones RAW
    \item Migración lógica filtros al servicio grabación
    \item Corrección problemas filtros canales auxiliares
    \item Migración polling → callbacks para evitar pérdida muestras
\end{itemize}

\section{Recursos y costes}

\subsection{Hardware}
\begin{itemize}
    \item PC de desarrollo: Intel i7-7700K + 32GB RAM
    \item Controlador PS4 DualShock para pruebas de entrada
    \item Casco EEG g.tec Unicorn para validación final (250Hz, 8 canales)
    \item Dispositivos EEG sintéticos BrainFlow para desarrollo (250Hz, múltiples canales)
\end{itemize}

\subsection{Software}
\begin{itemize}
    \item \textbf{Lenguajes}: Python 3.8+, JavaScript, HTML/CSS
    \item \textbf{Frameworks}: PyQt5, FastAPI, MNE-Python
    \item \textbf{Librerías principales}: BrainFlow, NumPy, SciPy, pynput, LSL
    \item \textbf{Herramientas}: Git, pre-commit, orjson, SDL2
    \item \textbf{Costes directos}: nulos (todas las licencias son libres)
\end{itemize}

\subsection{Estimación temporal}
\begin{itemize}
    \item \textbf{Desarrollo efectivo}: ~70 días (con pausa de 2 semanas)
    \item \textbf{Dedicación estimada}: ~350 horas
    \item \textbf{Commits realizados}: 200+ (desarrollo altamente iterativo)
\end{itemize}

\section{Análisis de riesgos y mitigación}

\begin{table}[hp!]
\centering
\resizebox{\textwidth}{!}{
\begin{tabular}{l|c|c|l}
\rowcolor{udcpink!25}
\textbf{Riesgo} & \textbf{Prob.} & \textbf{Impacto} & \textbf{Estrategia de mitigación} \\
\hline
Fallos repetidos integración módulos & Alta & Alta & Múltiples enfoques: LSL monolítico → SOLID → APIs \\
Escalabilidad código monolítico & Alta & Alta & Refactorización SOLID tras primer fallo integración \\
Dependencias GUI en modo headless & Media & Media & Identificadas pero no resueltas completamente \\
Rendimiento tiempo real crítico & Media & Alta & Eventos, callbacks, optimizaciones específicas \\
Curva aprendizaje tecnologías BCI & Media & Media & Investigación exhaustiva, sistema ayuda integrado \\
Sincronización temporal precisa & Alta & Media & Múltiples estrategias hasta solución viable \\
Dependencias bibliotecas externas & Baja & Alta & APIs estables, manejo robusto excepciones \\
\end{tabular}
}
\caption{Principales riesgos identificados y estrategias de mitigación}
\label{tab:riesgos}
\end{table}

\section{Seguimiento y control}

\subsection{Control de versiones}
\begin{itemize}
    \item \textbf{Repositorio Git}: más de 200 commits documentando progreso
    \item \textbf{Estrategia}: desarrollo directo en main con commits frecuentes
    \item \textbf{Versionado}: etiquetas semánticas para releases funcionales
\end{itemize}

\subsection{Metodología de desarrollo}
\begin{itemize}
    \item \textbf{Integración continua}: pre-commit hooks para calidad código
    \item \textbf{Type hints}: tipado progresivo, sesiones masivas tipado
    \item \textbf{Documentación}: generación automática documentación técnica
\end{itemize}

\subsection{Hitos principales}
\begin{itemize}
    \item \textbf{Milestone 1}: Módulos básicos + primer fallo integración (19/04)
    \item \textbf{Milestone 2}: Refactorización SOLID + segundo fallo (30/04)
    \item \textbf{Milestone 3}: Arquitectura API funcional (22/05)
    \item \textbf{Milestone 4}: Sistema completo con funcionalidades avanzadas (22/06)
\end{itemize}

\section{Desviaciones y lecciones aprendidas}

\subsection{Evolución metodológica}
\begin{itemize}
    \item \textbf{Julia → Python}: Precompilación y debugging complejo incompatibles con desarrollo iterativo
    \item \textbf{Dos fallos integración consecutivos}: Monolítica con LSL → SOLID con LSL → arquitectura API
    \item \textbf{tkinter → PyQt5}: Limitaciones rendimiento obligaron cambio framework UI
    \item \textbf{Polling → eventos → callbacks}: Evolución arquitectural por problemas rendimiento
\end{itemize}

\subsection{Problemas no resueltos}
\begin{itemize}
    \item \textbf{Modo headless CLI}: Dependencias Qt impiden funcionamiento completamente autónomo
    \item \textbf{Separación logs CLI}: Mezclado logs aplicación principal con CLI
    \item \textbf{Exportación EDF/CSV}: Fallos pausaron desarrollo módulo integración final
\end{itemize}

\subsection{Decisiones arquitecturales exitosas}
\begin{itemize}
    \item \textbf{APIs como backbone}: Solución robusta tras múltiples fallos integración directa
    \item \textbf{DataAccessService centralizado}: Eliminó duplicaciones y inconsistencias
    \item \textbf{Callbacks vs polling}: Solucionó pérdida muestras (10-30\%) en grabación
    \item \textbf{Arquitectura dirigida por eventos}: Mejoró rendimiento PS4 drasticamente
\end{itemize}